% ============================================================
%  sections/results.tex  –  Results
% ============================================================

\subsection{Baseline OLS estimates}

Table~\ref{tab:main-ols} reports the core log--log estimates linking transaction quantity to unit price.  The pooled specification implies an elasticity of $-0.551$, and the coefficient becomes slightly more negative as controls are added.  Once product quality and demographics are included, the estimate is $-0.633$.  We use the state-by-month fixed-effects specification in column~(6) as the main benchmark because it keeps the full analytic sample while absorbing state-specific monthly shocks, the richest time-varying market controls available in the main OLS block.  In that benchmark, the elasticity is $-0.652$.  The preceding state-only and state-plus-month specifications produce very similar elasticities of $-0.643$, so the estimate remains in a narrow band as the specification absorbs richer geographic and temporal heterogeneity.  This stability suggests that the baseline association is not driven solely by broad local or seasonal shocks.

The magnitude is also consistent with the transaction-based cannabis literature.  It is close to the crowdsourced estimate of \citet{davis2016crowd} for the United States and falls comfortably within the range of geography-based estimates summarized in Section~\ref{sec:litreview}, including \citet{halcoussis2017geo}.  In substantive terms, the benchmark coefficient implies that a 10\% higher unit price is associated with roughly a 6.5\% smaller transaction quantity. 

\begin{table}[H]
\centering
\caption{OLS estimates of the price elasticity of demand for marijuana}
\label{tab:main-ols}
\centering
\resizebox{\linewidth}{!}{%
\begin{tabular}[t]{lcccccc}
\toprule
  & (1) Pooled & (2) +Quality & (3) +Demog. & (4) +State FE & (5) +Month FE & (6) St×Mo FE\\
\midrule
ln(Price) & -0.551*** & -0.615*** & -0.633*** & -0.643*** & -0.643*** & -0.652***\\
 & (0.016) & (0.017) & (0.017) & (0.020) & (0.021) & (0.022)\\
Good quality &  & 0.651*** & 0.617*** & 0.600*** & 0.595*** & 0.603***\\
 &  & (0.044) & (0.044) & (0.060) & (0.059) & (0.059)\\
Age &  &  & 0.026*** & 0.024*** & 0.024*** & 0.022***\\
 &  &  & (0.003) & (0.004) & (0.004) & (0.005)\\
Gender: Other &  &  & 0.077 & 0.121 & 0.138 & 0.165\\
 &  &  & (0.171) & (0.158) & (0.159) & (0.163)\\
Gender: Women &  &  & 0.016 & 0.027 & 0.016 & 0.034\\
 &  &  & (0.052) & (0.054) & (0.053) & (0.059)\\
Education: Middle school &  &  & -0.002 & 0.043 & 0.022 & 0.010\\
 &  &  & (0.236) & (0.220) & (0.229) & (0.206)\\
Education: High school &  &  & 0.025 & 0.084 & 0.063 & 0.070\\
 &  &  & (0.217) & (0.201) & (0.210) & (0.180)\\
Education: Bachelor &  &  & 0.287 & 0.347 & 0.322 & 0.348*\\
 &  &  & (0.217) & (0.208) & (0.218) & (0.184)\\
Education: Graduate &  &  & 0.339 & 0.415* & 0.397 & 0.415*\\
 &  &  & (0.226) & (0.234) & (0.242) & (0.209)\\
\midrule
Num.Obs. & 3293 & 3293 & 3293 & 3293 & 3293 & 3293\\
R2 & 0.288 & 0.326 & 0.367 & 0.402 & 0.408 & 0.448\\
R2 Adj. & 0.288 & 0.326 & 0.365 & 0.395 & 0.399 & 0.401\\
FE: state\_code &  &  &  & X & X & \\
FE: month\_str &  &  &  &  & X & \\
FE: state\_month &  &  &  &  &  & X\\
\bottomrule
\multicolumn{7}{l}{\rule{0pt}{1em}* p $<$ 0.1, ** p $<$ 0.05, *** p $<$ 0.01}\\
\multicolumn{7}{l}{\rule{0pt}{1em}Dependent variable: ln(quantity in grams).}\\
\multicolumn{7}{l}{\rule{0pt}{1em}Robust SEs in parentheses (cols 1--3); clustered at state level (cols 4--6).}\\
\end{tabular}
}
\end{table}


\subsection{Quality heterogeneity}

Table~\ref{tab:quality-segmented} shows that the average elasticity masks substantial heterogeneity across product tiers.  When the sample is split by perceived quality, seedless marijuana has an estimated elasticity of $-0.613$, whereas seeded marijuana has an elasticity of $-0.891$.  The pooled interaction specification delivers the same message.  In that column, the coefficient on $\ln(\text{Price})$ is $-0.873$ for the reference category, and the interaction term of $0.259$ implies an elasticity of about $-0.614$ for good-quality transactions.  The close match between the split-sample and interaction estimates suggests that the quality gradient is not an artifact of changing samples.  All three columns retain the same state-by-month fixed effects and demographic controls as the benchmark, so the comparison isolates quality-tier heterogeneity rather than changes in the control set.

One interpretation is that seeded marijuana functions as a less differentiated lower-tier product, for which consumers may have closer substitutes across sellers.  Seedless marijuana, by contrast, may attract buyers with stronger preferences for quality or fewer close substitutes within the same tier.  Although the descriptive statistics show a large price gap between quality tiers, the slope estimates still point to greater price responsiveness in the lower-quality segment.

\begin{table}[H]
\centering
\caption{Quality-segmented demand estimates}
\label{tab:quality-segmented}
\centering
\begin{tabular}[t]{lccc}
\toprule
  & (7a) Good only & (7b) Bad only & (8) Interaction\\
\midrule
ln(Price) & -0.613*** & -0.891*** & -0.873***\\
 & (0.022) & (0.041) & (0.036)\\
Good quality &  &  & -0.019\\
 &  &  & (0.118)\\
ln(Price) x Good &  &  & 0.259***\\
 &  &  & (0.037)\\
Age & 0.023*** & 0.022*** & 0.023***\\
 & (0.005) & (0.007) & (0.005)\\
Gender: Other & 0.174 & 0.305 & 0.161\\
 & (0.208) & (0.184) & (0.163)\\
Gender: Women & -0.005 & 0.171 & 0.037\\
 & (0.059) & (0.113) & (0.059)\\
Education: Middle school & -0.124 & 0.193 & 0.017\\
 & (0.262) & (0.242) & (0.220)\\
Education: High school & -0.021 & 0.234* & 0.079\\
 & (0.270) & (0.116) & (0.196)\\
Education: Bachelor & 0.270 & 0.496*** & 0.353*\\
 & (0.263) & (0.057) & (0.198)\\
Education: Graduate & 0.345 & 0.539*** & 0.418*\\
 & (0.294) & (0.135) & (0.219)\\
\midrule
Num.Obs. & 2636 & 657 & 3293\\
R2 & 0.442 & 0.644 & 0.454\\
R2 Adj. & 0.385 & 0.514 & 0.408\\
\bottomrule
\multicolumn{4}{l}{\rule{0pt}{1em}* p $<$ 0.1, ** p $<$ 0.05, *** p $<$ 0.01}\\
\multicolumn{4}{l}{\rule{0pt}{1em}Dependent variable: ln(quantity in grams).}\\
\multicolumn{4}{l}{\rule{0pt}{1em}All columns include state-by-month FE and demographics. SEs clustered at state level.}\\
\end{tabular}
\end{table}


\subsection{Instrumental-variables diagnostics}

Table~\ref{tab:iv-combined} summarizes the IV diagnostics using the OSM/OSRM driving-time matrices and MUCD marijuana-seizure records described in Section~\ref{sec:data}.  The table follows the same order as the empirical strategy: OSM driving time from production hubs, lagged seizure-gravity indexes, and external Hausman-type instruments.  All columns include the benchmark state-by-month fixed effects and demographic controls; the external Hausman columns also include quality-by-month fixed effects, so the standalone quality indicator is absorbed in those specifications.  The first geography-based diagnostic replicates the production-hub distance logic used in prior cannabis studies with OSM driving time rather than straight-line distance.  It remains uninformative in this setting: the IV coefficient is $0.793$ with a standard error of $2.647$.  The one-month-lag hub-only marijuana-seizure gravity instrument yields $-0.653$, nearly identical to the benchmark OLS estimate, while the stricter outside-region seizure-gravity instrument yields $-0.680$ but with a larger standard error.  The two external Hausman instruments, which exclude the buyer's own state or macroregion when averaging prices within quality--month cells, yield negative elasticities of $-0.712$ and $-1.025$, respectively.  Because purchases in the estimation sample run from January through September~2019, the one-month-lag seizure-gravity specifications use MUCD marijuana seizures from December~2018 through August~2019 and can be merged to all 3{,}293 transactions.

\begin{table}[H]
\centering
\caption{OLS and IV estimates: alternative identification strategies}
\label{tab:iv-combined}
\centering
\resizebox{\linewidth}{!}{%
\begin{tabular}[t]{lcccc}
\toprule
  & (1) OLS & (2) IV: OSM driving time & (3) IV: Hausman & (4) IV: gravity\\
\midrule
ln(Price) & -0.652*** &  &  & \\
 & (0.022) &  &  & \\
ln(Price) [IV] &  & 0.793 & -0.648*** & -0.735***\\
 &  & (2.647) & (0.068) & (0.248)\\
Good quality & 0.603*** & -0.688 & 0.590*** & 0.677**\\
 & (0.059) & (2.363) & (0.069) & (0.250)\\
Age & 0.022*** & 0.025*** & 0.022*** & 0.022***\\
 & (0.005) & (0.009) & (0.005) & (0.005)\\
Gender: Other & 0.165 & 0.235 & 0.213 & 0.161\\
 & (0.163) & (0.345) & (0.178) & (0.165)\\
Gender: Women & 0.034 & 0.024 & 0.029 & 0.035\\
 & (0.059) & (0.135) & (0.059) & (0.057)\\
Education: Middle school & 0.010 & 0.485 & -0.013 & -0.018\\
 & (0.206) & (0.942) & (0.207) & (0.214)\\
Education: High school & 0.070 & 0.358 & 0.064 & 0.053\\
 & (0.180) & (0.587) & (0.180) & (0.177)\\
Education: Bachelor & 0.348* & 0.372 & 0.354* & 0.347*\\
 & (0.184) & (0.291) & (0.183) & (0.180)\\
Education: Graduate & 0.415* & 0.195 & 0.429** & 0.427*\\
 & (0.209) & (0.562) & (0.208) & (0.219)\\
\midrule
Num.Obs. & 3293 & 3293 & 3173 & 3293\\
R2 & 0.448 & 0.136 & 0.133 & 0.137\\
R2 Adj. & 0.401 & 0.063 & 0.075 & 0.064\\
\bottomrule
\multicolumn{5}{l}{\rule{0pt}{1em}* p $<$ 0.1, ** p $<$ 0.05, *** p $<$ 0.01}\\
\multicolumn{5}{l}{\rule{0pt}{1em}Dependent variable: ln(quantity in grams).}\\
\multicolumn{5}{l}{\rule{0pt}{1em}All columns include state-by-month FE and demographic controls.}\\
\multicolumn{5}{l}{\rule{0pt}{1em}OSM driving-time IV: instrument = ln(OSM driving time from nearest constructed production-hub point).}\\
\multicolumn{5}{l}{\rule{0pt}{1em}Hausman IV: instrument = leave-one-out mean ln(price) within state-quality-month cell.}\\
\multicolumn{5}{l}{\rule{0pt}{1em}Gravity IV: instrument = ln(1 + one-month-lag seizure-gravity index).}\\
\multicolumn{5}{l}{\rule{0pt}{1em}For weak geography-based IVs, inference should rely on Anderson-Rubin confidence sets rather than conventional significance stars.}\\
\multicolumn{5}{l}{\rule{0pt}{1em}SEs clustered at state level.}\\
\multicolumn{5}{l}{\rule{0pt}{1em}* p $<$ 0.1, ** p $<$ 0.05, *** p $<$ 0.01.}\\
\end{tabular}
}
\end{table}


Table~\ref{tab:iv-diagnostics} reports the corresponding first-stage and weak-instrument diagnostics, while repeating the IV coefficient from Table~\ref{tab:iv-combined} for readability.  The OSM production-hub instrument has a first-stage F-statistic of only $0.60$, well below conventional thresholds for reliable IV inference \citep{staiger1997instrumental,stock2005testing}.

The seizure-gravity diagnostics show a useful strength--exclusion tradeoff.  The hub-only lagged marijuana-seizure gravity instrument clears the conventional first-stage threshold ($F=15.64$), produces an IV estimate of $-0.653$, and has an Anderson--Rubin confidence set of $[-1.08,-0.41]$.  This is the closest IV estimate to the benchmark and is consistent with the interpretation that upstream enforcement shocks in production hubs shift retail prices.  However, the hub-only design still allows enforcement variation from broad production regions that may be connected to other market conditions.  The stricter outside-region seizure instrument addresses that concern more directly, but its first stage falls to $F=6.82$ and the Anderson--Rubin confidence set spans the full grid reported in the table.  For weak specifications, inference should be read through the Anderson--Rubin confidence sets rather than conventional second-stage significance stars \citep{anderson1949estimation}.

The external Hausman instruments are statistically stronger: excluding the buyer's own state gives $F=51.61$, and excluding the buyer's macroregion gives $F=21.73$.  The associated Anderson--Rubin intervals exclude zero, but the exclusion restriction remains an assumption about common cross-market price shocks rather than an empirical fact.  The Wu--Hausman p-values of $0.566$ and $0.152$ therefore indicate that OLS and these IV estimates are not statistically distinguishable under those maintained assumptions; they do not prove exogeneity.

\begin{table}[H]
\centering
\caption{Instrumental-variables estimates and diagnostics}
\label{tab:iv-diagnostics}
\centering
\resizebox{\linewidth}{!}{%
\begin{tabular}[t]{lccc}
\toprule
Diagnostic & OSM driving-time IV & Hausman IV & Seizure-gravity IV\\
\midrule
First-stage F & 0.60 & 169.09 & 7.07\\
IV estimate & 0.793 (2.647) & -0.648 (0.068) & -0.735 (0.248)\\
Wu--Hausman p-value & 0.335 & 0.951 & 0.740\\
Anderson--Rubin 95\% CI & [-2.50, 0.50] & [-0.79, -0.52] & [-2.50, 0.42]\\
\bottomrule
\multicolumn{4}{l}{\rule{0pt}{1em}All diagnostics use the state-by-month FE specifications in Table~\ref{tab:iv-combined}.}\\
\multicolumn{4}{l}{\rule{0pt}{1em}IV estimate entries reproduce the coefficient on ln(Price) [IV] and clustered SE from Table~\ref{tab:iv-combined}.}\\
\multicolumn{4}{l}{\rule{0pt}{1em}Seizure-gravity diagnostics use one-month-lagged MUCD marijuana seizures from December 2018--August 2019.}\\
\multicolumn{4}{l}{\rule{0pt}{1em}For weak geography-based IVs, inference should rely on Anderson-Rubin confidence sets rather than conventional significance stars.}\\
\multicolumn{4}{l}{\rule{0pt}{1em}OSM and seizure-gravity have first-stage F-statistics below 10; Hausman has a strong first stage but relies on a leave-one-out exclusion restriction.}\\
\end{tabular}
}
\end{table}


Taken together, the IV evidence supports a cautious reading of the benchmark.  The weak OSM result shows that production-hub driving time alone is not enough to identify a strong first stage in this sample.  The hub-only lagged seizure-gravity instrument reproduces the OLS magnitude with a first stage above 10, while the stricter outside-region version shows that stronger exclusion restrictions can come with weaker first-stage power.  The external Hausman instruments show that cross-market price variation produces negative elasticities of similar order.  We therefore treat the IV block as triangulation around the transaction-level price--quantity gradient, not as a standalone causal proof that supersedes the benchmark OLS specification.

\subsection{Robustness and sensitivity}

The baseline elasticity remains negative and economically meaningful under several robustness checks.  Table~\ref{tab:robustness} first shows that trimming the 1st and 99th percentiles of price and quantity moves the benchmark estimate only from $-0.652$ to $-0.620$.  Excluding round-number quantities, a direct response to heaping documented in the raw quantity distribution, yields an elasticity of $-0.576$.  Both checks therefore reduce the magnitude slightly, but neither eliminates the central negative price--quantity gradient.

\begin{table}[H]
\centering
\caption{Robustness: outlier trimming and heaping}
\label{tab:robustness}
\centering
\begin{tabular}[t]{lccc}
\toprule
  & (6) Benchmark & (9) Trimmed & (10) No heap\\
\midrule
ln(Price) & -0.652*** & -0.620*** & -0.576***\\
 & (0.022) & (0.022) & (0.028)\\
Good quality & 0.603*** & 0.586*** & 0.511***\\
 & (0.059) & (0.058) & (0.088)\\
Age & 0.022*** & 0.023*** & 0.009\\
 & (0.005) & (0.005) & (0.009)\\
Gender: Other & 0.165 & 0.071 & -0.250\\
 & (0.163) & (0.132) & (0.478)\\
Gender: Women & 0.034 & -0.012 & -0.002\\
 & (0.059) & (0.062) & (0.057)\\
Education: Middle school & 0.010 & -0.035 & -0.540***\\
 & (0.206) & (0.203) & (0.081)\\
Education: High school & 0.070 & 0.008 & -0.348***\\
 & (0.180) & (0.186) & (0.071)\\
Education: Bachelor & 0.348* & 0.271 & -0.170*\\
 & (0.184) & (0.195) & (0.087)\\
Education: Graduate & 0.415* & 0.313 & -0.002\\
 & (0.209) & (0.227) & (0.125)\\
\midrule
Num.Obs. & 3293 & 3198 & 1417\\
R2 & 0.448 & 0.406 & 0.401\\
R2 Adj. & 0.401 & 0.356 & 0.301\\
\bottomrule
\multicolumn{4}{l}{\rule{0pt}{1em}* p $<$ 0.1, ** p $<$ 0.05, *** p $<$ 0.01}\\
\multicolumn{4}{l}{\rule{0pt}{1em}Dependent variable: ln(quantity in grams).}\\
\multicolumn{4}{l}{\rule{0pt}{1em}All columns include state-by-month FE plus demographics. SEs clustered at state level.}\\
\multicolumn{4}{l}{\rule{0pt}{1em}Trimmed = drop 1st/99th percentile of price and quantity.}\\
\multicolumn{4}{l}{\rule{0pt}{1em}No heap = exclude round-number quantities (multiples of 5 or near 28g).}\\
\multicolumn{4}{l}{\rule{0pt}{1em}* p $<$ 0.1, ** p $<$ 0.05, *** p $<$ 0.01.}\\
\end{tabular}
\end{table}


Table~\ref{tab:muni-fe} provides an additional spatial credibility check by adding municipality fixed effects to the state-by-month benchmark.  This comparison is not purely a change in fixed effects, since the municipality-fixed-effects specification drops municipalities with one observation and uses $3{,}127$ observations rather than $3{,}293$.  Even in that restricted sample, the elasticity changes only from $-0.652$ to $-0.664$.  The augmented specification therefore absorbs both state-month shocks and time-invariant municipality heterogeneity, yet the estimated elasticity changes by only about one hundredth.

\begin{table}[H]
\centering
\caption{Demand estimates with alternative fixed-effects structures}
\label{tab:muni-fe}
\centering
\resizebox{\linewidth}{!}{%
\begin{tabular}[t]{lcc}
\toprule
  & (1) St×Mo FE & (2) St×Mo + Muni FE\\
\midrule
ln(Price) & -0.652*** & -0.664***\\
 & (0.022) & (0.021)\\
Good quality & 0.603*** & 0.571***\\
 & (0.059) & (0.053)\\
Age & 0.022*** & 0.022***\\
 & (0.005) & (0.004)\\
Gender: Other & 0.165 & 0.073\\
 & (0.163) & (0.143)\\
Gender: Women & 0.034 & -0.008\\
 & (0.059) & (0.060)\\
Education: Middle school & 0.010 & -0.215\\
 & (0.206) & (0.199)\\
Education: High school & 0.070 & -0.162\\
 & (0.180) & (0.194)\\
Education: Bachelor & 0.348* & 0.092\\
 & (0.184) & (0.198)\\
Education: Graduate & 0.415* & 0.175\\
 & (0.209) & (0.231)\\
\midrule
Num.Obs. & 3293 & 3127\\
R2 & 0.448 & 0.494\\
R2 Adj. & 0.401 & 0.413\\
\bottomrule
\multicolumn{3}{l}{\rule{0pt}{1em}* p $<$ 0.1, ** p $<$ 0.05, *** p $<$ 0.01}\\
\multicolumn{3}{l}{\rule{0pt}{1em}Dependent variable: ln(quantity in grams).}\\
\multicolumn{3}{l}{\rule{0pt}{1em}All columns include state-by-month FE and demographic controls.}\\
\multicolumn{3}{l}{\rule{0pt}{1em}Col 2 additionally includes municipality FE and restricts to municipalities with at least 2 observations.}\\
\multicolumn{3}{l}{\rule{0pt}{1em}SEs clustered at state level.}\\
\end{tabular}
}
\end{table}


With only 32 state clusters, we also complement conventional clustered standard errors with a wild cluster bootstrap in the sense of \citet{cameron2008bootstrap} and \citet{roodman2019fast}.  Table~\ref{tab:bootstrap-inference} reports the conventional state-clustered interval, the wild cluster bootstrap interval using Webb weights and 9{,}999 replications, and a pairs cluster bootstrap sensitivity check with the same number of replications.  The wild-bootstrap interval remains tightly centered around the benchmark elasticity and excludes zero by a wide margin, so the main inference does not appear to be an artifact of conventional clustered standard errors.

\begin{table}[H]
\centering
\caption{Bootstrap inference for the benchmark elasticity}
\label{tab:bootstrap-inference}
\centering
\begin{tabular}[t]{lccc}
\toprule
Method & Estimate & SE & 95\% CI\\
\midrule
State-clustered SE & -0.652 & 0.023 & [-0.696, -0.607]\\
Wild cluster bootstrap & -0.652 & -- & [-0.715, -0.602]\\
Pairs cluster bootstrap & -0.653 & 0.023 & [-0.704, -0.610]\\
\bottomrule
\multicolumn{4}{l}{\rule{0pt}{1em}Benchmark model: state-by-month FE, quality, age, gender, and education controls.}\\
\multicolumn{4}{l}{\rule{0pt}{1em}Wild bootstrap uses Webb weights, 9,999 replications, and state clusters.}\\
\multicolumn{4}{l}{\rule{0pt}{1em}Pairs cluster bootstrap resamples states with replacement across 1,000 replications.}\\
\end{tabular}
\end{table}


Finally, Table~\ref{tab:oster-bounds} reports the coefficient-stability exercise of \citet{oster2019unobservable}.  The calculation uses the short and long specifications from columns~(1) and~(6) of Table~\ref{tab:main-ols}, with $R_{\max}$ set to $1.3R^2_{\text{long}}$.  Because adding observables moves the elasticity away from zero, the raw $\delta^{*}$ needed to set the coefficient to zero is negative; the table reports its absolute value, $|\delta^{*}|=7.70$.  Thus, explaining away the benchmark would require unobserved selection that is large and works in the opposite direction from the observed controls.  This is not a proof of identification, but it makes an omitted-variable explanation of the benchmark elasticity less plausible.

\begin{table}[H]
\centering
\caption{Oster coefficient-stability diagnostics}
\label{tab:oster-bounds}
\centering
\begin{tabular}[t]{p{0.62\linewidth}r}
\toprule
Statistic & Estimate\\
\midrule
Short-specification elasticity & -0.551\\
Short-specification $R^2$ & 0.288\\
Long-specification elasticity & -0.652\\
Long-specification $R^2$ & 0.448\\
$R_{\max}=1.3R^2_{long}$ & 0.582\\
Bias-adjusted elasticity ($\delta=1$) & -0.736\\
$|\delta^*|$ to set elasticity to zero & 7.70\\
\bottomrule
\multicolumn{2}{p{0.78\linewidth}}{\footnotesize \rule{0pt}{1em}Short specification includes only ln(price); long specification adds controls plus state-by-month FE.}\\
\multicolumn{2}{p{0.78\linewidth}}{\footnotesize \rule{0pt}{1em}The table reports the absolute value of $\delta^*$; the raw calculation is negative because coefficients are negative.}\\
\end{tabular}
\end{table}

